% ─── AUTO-GERADO por generate_dataset.py (semente 42): 248 consultas ───

\section{Consultas simples geradas por \textit{templates} (G-S)}

Total desta fam\'ilia: \textbf{61 consultas}, listadas integralmente abaixo (a fun\c{c}\~ao geradora e o modelo de frase de cada consulta constam de \texttt{dataset.json}).

{\footnotesize
\begin{longtable}{|p{1.1cm}|p{1.3cm}|p{5.4cm}|p{6.3cm}|}
\hline
\textbf{ID} & \textbf{Cat.} & \textbf{Consulta} & \textbf{Leituras aceitas} \\ \hline
\endhead
GS001 & S & cartas do Acre & state: "Acre" \\ \hline
GS002 & S · A & cartas de AL & state: "Alagoas" \\ \hline
GS003 & S · A & cartas do AP & state: "Amapá" \\ \hline
GS004 & S & cartas do Amazonas & state: "Amazonas" \\ \hline
GS005 & S · A & cartas da bahia & state: "Bahia" \\ \hline
GS006 & S · A & cartas do CE & state: "Ceará" \\ \hline
GS007 & S · A & cartas do DF & state: "Distrito Federal" \\ \hline
GS008 & S · A & cartas do ES & state: "Espírito Santo" \\ \hline
GS009 & S & cartas de Goiás & state: "Goiás" \\ \hline
GS010 & S · A & cartas do MA & state: "Maranhão" \\ \hline
GS011 & S & cartas de Mato Grosso & state: "Mato Grosso" \\ \hline
GS012 & S & cartas de Mato Grosso do Sul & state: "Mato Grosso do Sul" \\ \hline
GS013 & S · A & cartas de MG & state: "Minas Gerais" \\ \hline
GS014 & S & cartas do Pará & state: "Pará" \\ \hline
GS015 & S · A & cartas da paraíba & state: "Paraíba" \\ \hline
GS016 & S · A & produtos na escala 1:1000 & scale: "1:1.000" \\ \hline
GS017 & S · A & produtos na escala 1:2000 & scale: "1:2.000" \\ \hline
GS018 & S · A & produtos na escala 1:5000 & scale: "1:5.000" \\ \hline
GS019 & S · A & produtos na escala 1:10000 & scale: "1:10.000" \\ \hline
GS020 & S · A & produtos na escala 1:25000 & scale: "1:25.000" \\ \hline
GS021 & S · A & produtos na escala de 50 mil & scale: "1:50.000" \\ \hline
GS022 & S · A & produtos na escala de 100 mil & scale: "1:100.000" \\ \hline
GS023 & S · A & produtos em 250k & scale: "1:250.000" \\ \hline
GS024 & S · A & cartas topográficas disponíveis & productType: "SCN Carta Topográfica Matricial" \\ \hline
GS025 & S · A & cartas vetoriais disponíveis & productType: "SCN Carta Topográfica Vetorial" \\ \hline
GS026 & S · A & ortoimagens disponíveis & productType: "SCN Carta Ortoimagem" \\ \hline
GS027 & S · A & ortoimagens banda P HH disponíveis & productType: "SCN Carta Ortoimagem Banda P Pol HH" \\ \hline
GS028 & S · A & ortoimagens banda X HH disponíveis & productType: "SCN Carta Ortoimagem Banda X Pol HH" \\ \hline
GS029 & S · A & MDTs disponíveis & productType: "MDT — RAM" \\ \hline
GS030 & S · A & MDSs disponíveis & productType: "MDS — RAM" \\ \hline
GS031 & S · A & cartas CIRC disponíveis & productType: "CIRC" \\ \hline
GS032 & S · A & mapas temáticos disponíveis & productType: "Cartas Temáticas Não SCN" \\ \hline
GS033 & S · A & produtos do primeiro cgeo & supplyArea: "1$^\circ$ Centro de Geoinformação" \\ \hline
GS034 & S · A & produtos do 2o cgeo & supplyArea: "2$^\circ$ Centro de Geoinformação" \\ \hline
GS035 & S · A & produtos do 3o cgeo & supplyArea: "3$^\circ$ Centro de Geoinformação" \\ \hline
GS036 & S · A & produtos do quarto cgeo & supplyArea: "4$^\circ$ Centro de Geoinformação" \\ \hline
GS037 & S · A & produtos do 5º CGEO & supplyArea: "5$^\circ$ Centro de Geoinformação" \\ \hline
GS038 & S & cartas de Brasília & city: "Brasília" \\ \hline
GS039 & S & cartas de Manaus & city: "Manaus" \\ \hline
GS040 & S & cartas de Porto Velho & city: "Porto Velho" \\ \hline
GS041 & S & cartas de Cuiabá & city: "Cuiabá" \\ \hline
GS042 & S & cartas de Fortaleza & city: "Fortaleza" \\ \hline
GS043 & S & cartas de Salvador & city: "Salvador" \\ \hline
GS044 & S & cartas de Belo Horizonte & city: "Belo Horizonte" \\ \hline
GS045 & S & cartas de Curitiba & city: "Curitiba" \\ \hline
GS046 & S & cartas de Porto Alegre & city: "Porto Alegre" \\ \hline
GS047 & S & cartas de Recife & city: "Recife" \\ \hline
GS048 & S · M & carta MI 2965-2-NE & keyword: "2965-2-NE" \\ \hline
GS049 & S · M & carta MI 2866-3 & keyword: "2866-3" \\ \hline
GS050 & S · M & carta MI 2901 & keyword: "2901" \\ \hline
GS051 & S · M & carta MI 530 & keyword: "530" \\ \hline
GS052 & S · M & carta MI 3010 & keyword: "3010" \\ \hline
GS053 & S · M & carta MI 2965 & keyword: "2965" \\ \hline
GS054 & S · M & carta MI 2866-2-SO & keyword: "2866-2-SO" \\ \hline
GS055 & S · M & carta MI 2901-4-NE & keyword: "2901-4-NE" \\ \hline
GS056 & S · M & folha INOM SF-22-Y-D & keyword: "SF-22-Y-D" \\ \hline
GS057 & S · M & folha INOM SF-22-Y-D-II & keyword: "SF-22-Y-D-II" \\ \hline
GS058 & S · M & folha INOM SF-22-Y-D-II-4 & keyword: "SF-22-Y-D-II-4" \\ \hline
GS059 & S · M & folha INOM SF-22-Y-D-II-4-SE & keyword: "SF-22-Y-D-II-4-SE" \\ \hline
GS060 & S · M & folha INOM SG-22-X-A & keyword: "SG-22-X-A" \\ \hline
GS061 & S · M & folha INOM SF-23-V-C & keyword: "SF-23-V-C" \\ \hline
\end{longtable}
\addtocounter{table}{-1}
}

\section{Consultas compostas geradas por \textit{templates} (G-C)}

Total desta fam\'ilia: \textbf{55 consultas}, listadas integralmente abaixo (a fun\c{c}\~ao geradora e o modelo de frase de cada consulta constam de \texttt{dataset.json}).

{\footnotesize
\begin{longtable}{|p{1.1cm}|p{1.3cm}|p{5.4cm}|p{6.3cm}|}
\hline
\textbf{ID} & \textbf{Cat.} & \textbf{Consulta} & \textbf{Leituras aceitas} \\ \hline
\endhead
GC001 & C · A & cartas do amapá na escala 1:2.000 & state: "Amapá"; scale: "1:2.000" \\ \hline
GC002 & C · A & cartas de Mato Grosso do Sul em 25k & state: "Mato Grosso do Sul"; scale: "1:25.000" \\ \hline
GC003 & C · A & cartas de santa catarina na escala 1:2.000 & state: "Santa Catarina"; scale: "1:2.000" \\ \hline
GC004 & C · A & cartas do Amapá na escala de 25 mil & state: "Amapá"; scale: "1:25.000" \\ \hline
GC005 & C · A & cartas do RJ na escala 1:2000 & state: "Rio de Janeiro"; scale: "1:2.000" \\ \hline
GC006 & C · A & cartas de RO em 25k & state: "Rondônia"; scale: "1:25.000" \\ \hline
GC007 & C · A & cartas do ES na escala 1:100000 & state: "Espírito Santo"; scale: "1:100.000" \\ \hline
GC008 & C · A & cartas da Paraíba em 50k & state: "Paraíba"; scale: "1:50.000" \\ \hline
GC009 & C · A & cartas de mato grosso do sul na escala de 10 mil & state: "Mato Grosso do Sul"; scale: "1:10.000" \\ \hline
GC010 & C · A & cartas de Goiás na escala 1:2000 & state: "Goiás"; scale: "1:2.000" \\ \hline
GC011 & C · A & cartas do Piauí na escala 1:10000 & state: "Piauí"; scale: "1:10.000" \\ \hline
GC012 & C · A & cartas da paraíba na escala 1:25000 & state: "Paraíba"; scale: "1:25.000" \\ \hline
GC013 & C · A & cartas de rondônia na escala 1:1000 & state: "Rondônia"; scale: "1:1.000" \\ \hline
GC014 & C · A & cartas do TO na escala 1:50000 & state: "Tocantins"; scale: "1:50.000" \\ \hline
GC015 & C · A & cartas de GO na escala de 10 mil & state: "Goiás"; scale: "1:10.000" \\ \hline
GC016 & C · A & cartas de roraima na escala de 10 mil & state: "Roraima"; scale: "1:10.000" \\ \hline
GC017 & C · A & cartas do paraná na escala 1:250000 & state: "Paraná"; scale: "1:250.000" \\ \hline
GC018 & C · A & cartas do Piauí em escala detalhada & state: "Piauí"; scale: "1:25.000" \\ \hline
GC019 & C · A & cartas do rio de janeiro em 50k & state: "Rio de Janeiro"; scale: "1:50.000"; aceita-se também city no lugar de state \\ \hline
GC020 & C · A & cartas da Bahia em 250k & state: "Bahia"; scale: "1:250.000" \\ \hline
GC021 & C · A & cartas de SP na escala 1:2000 & state: "São Paulo"; scale: "1:2.000" \\ \hline
GC022 & C · A & cartas do RS na escala de 100 mil & state: "Rio Grande do Sul"; scale: "1:100.000" \\ \hline
GC023 & C · A & cartas do amapá na escala de 100 mil & state: "Amapá"; scale: "1:100.000" \\ \hline
GC024 & C · A & cartas da Paraíba em 25k & state: "Paraíba"; scale: "1:25.000" \\ \hline
GC025 & C & cartas de Rondônia na escala 1:2.000 & state: "Rondônia"; scale: "1:2.000" \\ \hline
GC026 & C · A & cartas de São Paulo em 50k & state: "São Paulo"; scale: "1:50.000"; aceita-se também city no lugar de state \\ \hline
GC027 & C · A & cartas do maranhão na escala 1:5.000 & state: "Maranhão"; scale: "1:5.000" \\ \hline
GC028 & C · A & cartas do Acre na escala 1:25000 & state: "Acre"; scale: "1:25.000" \\ \hline
GC029 & C · A & cartas de PE na escala 1:25000 & state: "Pernambuco"; scale: "1:25.000" \\ \hline
GC030 & C · A & cartas da bahia na escala 1:5000 & state: "Bahia"; scale: "1:5.000" \\ \hline
GC031 & C · A & mapas de Recife na escala 1:50000 & city: "Recife"; scale: "1:50.000" \\ \hline
GC032 & C & mapas de Brasília na escala 1:2.000 & city: "Brasília"; scale: "1:2.000" \\ \hline
GC033 & C · A & mapas de Fortaleza na escala 1:10000 & city: "Fortaleza"; scale: "1:10.000" \\ \hline
GC034 & C · A & mapas de Cuiabá na escala 1:2000 & city: "Cuiabá"; scale: "1:2.000" \\ \hline
GC035 & C · A & mapas de Belém em 250k & city: "Belém"; scale: "1:250.000" \\ \hline
GC036 & C · A & mapas de Porto Alegre na escala 1:5000 & city: "Porto Alegre"; scale: "1:5.000" \\ \hline
GC037 & C · A & mapas de Goiânia na escala 1:250000 & city: "Goiânia"; scale: "1:250.000" \\ \hline
GC038 & C · A & mapas de Fortaleza em 100k & city: "Fortaleza"; scale: "1:100.000" \\ \hline
GC039 & C · A & cartas temáticas de São Paulo & productType: "Cartas Temáticas Não SCN"; state: "São Paulo"; aceita-se também city no lugar de state \\ \hline
GC040 & C · A & ortoimagens banda X HH de minas & productType: "SCN Carta Ortoimagem Banda X Pol HH"; state: "Minas Gerais" \\ \hline
GC041 & C · A & cartas temáticas da PB & productType: "Cartas Temáticas Não SCN"; state: "Paraíba" \\ \hline
GC042 & C · A & ortoimagens banda P HH do amapá & productType: "SCN Carta Ortoimagem Banda P Pol HH"; state: "Amapá" \\ \hline
GC043 & C · A & cartas temáticas do Espírito Santo & productType: "Cartas Temáticas Não SCN"; state: "Espírito Santo" \\ \hline
GC044 & C · A & cartas topográficas do Amapá & productType: "SCN Carta Topográfica Matricial"; state: "Amapá" \\ \hline
GC045 & C · A & produtos do projeto copa das confederacoes no AP & project: "Copa das Confederações"; state: "Amapá" \\ \hline
GC046 & C · A & produtos do projeto olimpiadas em PE & project: "Olimpíadas Rio 2016"; state: "Pernambuco" \\ \hline
GC047 & C & produtos do projeto NGA-BECA no Distrito Federal & project: "NGA-BECA"; state: "Distrito Federal" \\ \hline
GC048 & C · A & produtos do projeto beca no espírito santo & project: "NGA-BECA"; state: "Espírito Santo" \\ \hline
GC049 & C · A & produtos do projeto copa das confederacoes no AM & project: "Copa das Confederações"; state: "Amazonas" \\ \hline
GC050 & C · A & produtos do projeto BCD de Rondônia no pará & project: "Base Cartográfica Digital de Rondônia"; state: "Pará" \\ \hline
GC051 & C · A & cartas do 1º CGEO na escala 1:2.000 & supplyArea: "1$^\circ$ Centro de Geoinformação"; scale: "1:2.000" \\ \hline
GC052 & C · A & cartas do 3o cgeo na escala 1:2000 & supplyArea: "3$^\circ$ Centro de Geoinformação"; scale: "1:2.000" \\ \hline
GC053 & C · A & cartas do 2o cgeo em pequena escala & supplyArea: "2$^\circ$ Centro de Geoinformação"; scale: "1:250.000" \\ \hline
GC054 & C · A & cartas do 2 CGEO na escala 1:25000 & supplyArea: "2$^\circ$ Centro de Geoinformação"; scale: "1:25.000" \\ \hline
GC055 & C · A & cartas do 1º CGEO em 250k & supplyArea: "1$^\circ$ Centro de Geoinformação"; scale: "1:250.000" \\ \hline
\end{longtable}
\addtocounter{table}{-1}
}

\section{Consultas com c\'odigo MI/INOM geradas por \textit{templates} (G-M)}

Total desta fam\'ilia: \textbf{28 consultas}, listadas integralmente abaixo (a fun\c{c}\~ao geradora e o modelo de frase de cada consulta constam de \texttt{dataset.json}).

{\footnotesize
\begin{longtable}{|p{1.1cm}|p{1.3cm}|p{5.4cm}|p{6.3cm}|}
\hline
\textbf{ID} & \textbf{Cat.} & \textbf{Consulta} & \textbf{Leituras aceitas} \\ \hline
\endhead
GM001 & C · M & folha MI 2965-2-NE do Rio Grande do Sul & keyword: "2965-2-NE"; state: "Rio Grande do Sul" \\ \hline
GM002 & C · M · A & folha MI 2866-3 do TO & keyword: "2866-3"; state: "Tocantins" \\ \hline
GM003 & C · M · A & folha MI 2901 do AP & keyword: "2901"; state: "Amapá" \\ \hline
GM004 & C · M · A & folha MI 530 do ceará & keyword: "530"; state: "Ceará" \\ \hline
GM005 & C · M · A & folha MI 3010 do paraná & keyword: "3010"; state: "Paraná" \\ \hline
GM006 & C · M · A & folha MI 2965 do distrito federal & keyword: "2965"; state: "Distrito Federal" \\ \hline
GM007 & C · M · A & folha MI 2866-2-SO de AL & keyword: "2866-2-SO"; state: "Alagoas" \\ \hline
GM008 & C · M · A & folha MI 2901-4-NE de MG & keyword: "2901-4-NE"; state: "Minas Gerais" \\ \hline
GM009 & C · M · A & folha MI 531-1-NE de minas & keyword: "531-1-NE"; state: "Minas Gerais" \\ \hline
GM010 & C · M & folha MI 2965-3 de Sergipe & keyword: "2965-3"; state: "Sergipe" \\ \hline
GM011 & C · M · A & folha MI 2866-3-NO do maranhão & keyword: "2866-3-NO"; state: "Maranhão" \\ \hline
GM012 & C · M & folha MI 3010-2-SE de Roraima & keyword: "3010-2-SE"; state: "Roraima" \\ \hline
GM013 & C · M · A & INOM SF-22-Y-D do tipo temática & keyword: "SF-22-Y-D"; productType: "Cartas Temáticas Não SCN" \\ \hline
GM014 & C · M · A & INOM SF-22-Y-D-II do tipo ortoimagem & keyword: "SF-22-Y-D-II"; productType: "SCN Carta Ortoimagem" \\ \hline
GM015 & C · M · A & INOM SF-22-Y-D-II-4 do tipo banda X HH & keyword: "SF-22-Y-D-II-4"; productType: "SCN Carta Ortoimagem Banda X Pol HH" \\ \hline
GM016 & C · M · A & INOM SF-22-Y-D-II-4-SE do tipo topográfica & keyword: "SF-22-Y-D-II-4-SE"; productType: "SCN Carta Topográfica Matricial" \\ \hline
GM017 & C · M · A & INOM SG-22-X-A do tipo MDT & keyword: "SG-22-X-A"; productType: "MDT — RAM" \\ \hline
GM018 & C · M · A & INOM SF-23-V-C do tipo topo & keyword: "SF-23-V-C"; productType: "SCN Carta Topográfica Matricial" \\ \hline
GM019 & C · M · A & INOM SG-21-Z-B do tipo temática & keyword: "SG-21-Z-B"; productType: "Cartas Temáticas Não SCN" \\ \hline
GM020 & C · M · A & INOM SH-22-W-A-III-1-NE do tipo topográfica & keyword: "SH-22-W-A-III-1-NE"; productType: "SCN Carta Topográfica Matricial" \\ \hline
GM021 & C · M · A & INOM SG-22-Z-D-II-3 do tipo ortoimagem & keyword: "SG-22-Z-D-II-3"; productType: "SCN Carta Ortoimagem" \\ \hline
GM022 & C · M · A & INOM SF-24-Y-A do tipo vetorial & keyword: "SF-24-Y-A"; productType: "SCN Carta Topográfica Vetorial" \\ \hline
GM023 & C · M · A & MI 2866-2-SO produzido pelo 1 CGEO & keyword: "2866-2-SO"; supplyArea: "1$^\circ$ Centro de Geoinformação" \\ \hline
GM024 & C · M · A & MI 2866-3 produzido pelo quinto cgeo & keyword: "2866-3"; supplyArea: "5$^\circ$ Centro de Geoinformação" \\ \hline
GM025 & C · M · A & MI 2965-3 produzido pelo 3o cgeo & keyword: "2965-3"; supplyArea: "3$^\circ$ Centro de Geoinformação" \\ \hline
GM026 & C · M · A & MI 530 produzido pelo 3o cgeo & keyword: "530"; supplyArea: "3$^\circ$ Centro de Geoinformação" \\ \hline
GM027 & C · M · A & MI 2965-2-NE produzido pelo 3 CGEO & keyword: "2965-2-NE"; supplyArea: "3$^\circ$ Centro de Geoinformação" \\ \hline
GM028 & C · M · A & MI 3010 produzido pelo segundo cgeo & keyword: "3010"; supplyArea: "2$^\circ$ Centro de Geoinformação" \\ \hline
\end{longtable}
\addtocounter{table}{-1}
}

\section{Consultas com tempo relativo geradas por \textit{templates} (G-T)}

Total desta fam\'ilia: \textbf{41 consultas}, listadas integralmente abaixo (a fun\c{c}\~ao geradora e o modelo de frase de cada consulta constam de \texttt{dataset.json}).

{\footnotesize
\begin{longtable}{|p{1.1cm}|p{1.3cm}|p{5.4cm}|p{6.3cm}|}
\hline
\textbf{ID} & \textbf{Cat.} & \textbf{Consulta} & \textbf{Leituras aceitas} \\ \hline
\endhead
GT001 & T & cartas publicadas esse ano & publicationPeriod: ano corrente \\ \hline
GT002 & T & cartas publicadas no ano passado & publicationPeriod: ano anterior, inteiro \\ \hline
GT003 & T & cartas publicadas no mês passado & publicationPeriod: mês anterior, inteiro \\ \hline
GT004 & T & cartas publicadas na semana passada & publicationPeriod: semana passada \\ \hline
GT005 & T & cartas publicadas nos últimos 3 meses & publicationPeriod: últimos 3 meses \\ \hline
GT006 & T & cartas publicadas nos últimos 6 meses & publicationPeriod: últimos 6 meses \\ \hline
GT007 & T & cartas publicadas nos últimos 5 anos & publicationPeriod: últimos 5 anos \\ \hline
GT008 & T & cartas publicadas desde 2020 & publicationPeriod: desde 2020 \\ \hline
GT009 & T & cartas publicadas depois de 2022 & publicationPeriod: depois de 2022 \\ \hline
GT010 & T & cartas publicadas antes de 2020 & publicationPeriod: antes de 2020 \\ \hline
GT011 & T & cartas publicadas entre 2022 e 2023 & publicationPeriod: 2022-01-01 a 2023-12-31 \\ \hline
GT012 & T & cartas publicadas no primeiro trimestre deste ano & publicationPeriod: 1º trimestre do ano corrente \\ \hline
GT013 & T & cartas publicadas no segundo semestre do ano passado & publicationPeriod: 2º semestre do ano anterior \\ \hline
GT014 & T & cartas publicadas nesta semana & publicationPeriod: semana corrente \\ \hline
GT015 & T & cartas publicadas hoje & publicationPeriod: dia da execução \\ \hline
GT016 & C · T · A & cartas de pernambuco publicadas desde 2020 & state: "Pernambuco"; publicationPeriod: desde 2020 \\ \hline
GT017 & C · T · A & cartas da bahia publicadas nos últimos 6 meses & state: "Bahia"; publicationPeriod: últimos 6 meses \\ \hline
GT018 & C · T · A & cartas do AP produzidas no segundo semestre do ano passado & state: "Amapá"; creationPeriod: 2º semestre do ano anterior \\ \hline
GT019 & C · T · A & cartas do MA publicadas no ano passado & state: "Maranhão"; publicationPeriod: ano anterior, inteiro \\ \hline
GT020 & C · T & cartas da Paraíba publicadas esse ano & state: "Paraíba"; publicationPeriod: ano corrente \\ \hline
GT021 & C · T · A & cartas de roraima elaboradas hoje & state: "Roraima"; creationPeriod: dia da execução \\ \hline
GT022 & C · T · A & cartas de sergipe publicadas no primeiro trimestre deste ano & state: "Sergipe"; publicationPeriod: 1º trimestre do ano corrente \\ \hline
GT023 & C · T & cartas de Rondônia publicadas depois de 2022 & state: "Rondônia"; publicationPeriod: depois de 2022 \\ \hline
GT024 & C · T & cartas do Maranhão criadas nos últimos 3 meses & state: "Maranhão"; creationPeriod: últimos 3 meses \\ \hline
GT025 & C · T · A & cartas da bahia publicadas antes de 2020 & state: "Bahia"; publicationPeriod: antes de 2020 \\ \hline
GT026 & C · T · A & produtos do 3o cgeo publicados no ano passado & supplyArea: "3$^\circ$ Centro de Geoinformação"; publicationPeriod: ano anterior, inteiro \\ \hline
GT027 & C · T · A & produtos do 3 CGEO publicados hoje & supplyArea: "3$^\circ$ Centro de Geoinformação"; publicationPeriod: dia da execução \\ \hline
GT028 & C · T · A & produtos do 3º CGEO criados no primeiro trimestre deste ano & supplyArea: "3$^\circ$ Centro de Geoinformação"; creationPeriod: 1º trimestre do ano corrente \\ \hline
GT029 & C · T · A & produtos do quarto cgeo publicados depois de 2022 & supplyArea: "4$^\circ$ Centro de Geoinformação"; publicationPeriod: depois de 2022 \\ \hline
GT030 & C · T · A & produtos do 1º CGEO publicados no mês passado & supplyArea: "1$^\circ$ Centro de Geoinformação"; publicationPeriod: mês anterior, inteiro \\ \hline
GT031 & C · T · A & produtos do terceiro cgeo criados nos últimos 3 meses & supplyArea: "3$^\circ$ Centro de Geoinformação"; creationPeriod: últimos 3 meses \\ \hline
GT032 & C · T · A & produtos do 2 CGEO publicados antes de 2020 & supplyArea: "2$^\circ$ Centro de Geoinformação"; publicationPeriod: antes de 2020 \\ \hline
GT033 & C · T · A & produtos do 5 CGEO publicados na semana passada & supplyArea: "5$^\circ$ Centro de Geoinformação"; publicationPeriod: semana passada \\ \hline
GT034 & C · T & cartas na escala 1:5.000 publicadas depois de 2022 & scale: "1:5.000"; publicationPeriod: depois de 2022 \\ \hline
GT035 & C · T · A & cartas na escala 1:5000 publicadas esse ano & scale: "1:5.000"; publicationPeriod: ano corrente \\ \hline
GT036 & C · T · A & cartas na escala de 25 mil criadas no ano passado & scale: "1:25.000"; creationPeriod: ano anterior, inteiro \\ \hline
GT037 & C · T · A & cartas em 50k publicadas no segundo semestre do ano passado & scale: "1:50.000"; publicationPeriod: 2º semestre do ano anterior \\ \hline
GT038 & C · T · A & cartas na escala de 10 mil publicadas no mês passado & scale: "1:10.000"; publicationPeriod: mês anterior, inteiro \\ \hline
GT039 & C · T & cartas na escala 1:2.000 elaboradas hoje & scale: "1:2.000"; creationPeriod: dia da execução \\ \hline
GT040 & C · T · A & cartas na escala de 100 mil publicadas nesta semana & scale: "1:100.000"; publicationPeriod: semana corrente \\ \hline
GT041 & C · T · A & cartas na escala 1:5000 publicadas nos últimos 6 meses & scale: "1:5.000"; publicationPeriod: últimos 6 meses \\ \hline
\end{longtable}
\addtocounter{table}{-1}
}

\section{Consultas com ordena\c{c}\~ao geradas por \textit{templates} (G-O)}

Total desta fam\'ilia: \textbf{19 consultas}, listadas integralmente abaixo (a fun\c{c}\~ao geradora e o modelo de frase de cada consulta constam de \texttt{dataset.json}).

{\footnotesize
\begin{longtable}{|p{1.1cm}|p{1.3cm}|p{5.4cm}|p{6.3cm}|}
\hline
\textbf{ID} & \textbf{Cat.} & \textbf{Consulta} & \textbf{Leituras aceitas} \\ \hline
\endhead
GO001 & C · O · A & a carta mais recente do CE & state: "Ceará"; sortField: "publicationDate" ou "creationDate"; sortDirection: "DESC"; limit: 1 (opcional) \\ \hline
GO002 & C · O · A & a carta mais antiga do PA & state: "Pará"; sortField: "publicationDate" ou "creationDate"; sortDirection: "ASC"; limit: 1 (opcional) \\ \hline
GO003 & C · O & as cartas mais recentes do Ceará & state: "Ceará"; sortField: "publicationDate" ou "creationDate"; sortDirection: "DESC" \\ \hline
GO004 & C · O · A & as cartas mais antigas de mato grosso & state: "Mato Grosso"; sortField: "publicationDate" ou "creationDate"; sortDirection: "ASC" \\ \hline
GO005 & C · O · A & a primeira carta produzida de sergipe & state: "Sergipe"; sortField: "creationDate"; sortDirection: "ASC"; limit: 1 (opcional) \\ \hline
GO006 & C · O · A & a última carta publicada de GO & state: "Goiás"; sortField: "publicationDate"; sortDirection: "DESC"; limit: 1 (opcional) \\ \hline
GO007 & C · O · A & as 3 cartas mais antigas de sergipe & state: "Sergipe"; sortField: "publicationDate" ou "creationDate"; sortDirection: "ASC"; limit: 3 \\ \hline
GO008 & C · O · A & as 5 publicações mais recentes de MG & state: "Minas Gerais"; sortField: "publicationDate"; sortDirection: "DESC"; limit: 5 \\ \hline
GO009 & C · O · A & as 10 cartas criadas mais recentemente do PR & state: "Paraná"; sortField: "creationDate"; sortDirection: "DESC"; limit: 10 \\ \hline
GO010 & C · O · A & cartas do distrito federal em ordem cronológica de publicação & state: "Distrito Federal"; sortField: "publicationDate"; sortDirection: "ASC" \\ \hline
GO011 & C · O & as cartas temáticas mais recentes & productType: "Cartas Temáticas Não SCN"; sortField: "publicationDate" ou "creationDate"; sortDirection: "DESC" \\ \hline
GO012 & C · O & as 3 cartas CIRC mais antigas & productType: "CIRC"; sortField: "publicationDate" ou "creationDate"; sortDirection: "ASC"; limit: 3 \\ \hline
GO013 & C · O & a última ortoimagem publicada & productType: "SCN Carta Ortoimagem"; sortField: "publicationDate"; sortDirection: "DESC"; limit: 1 (opcional) \\ \hline
GO014 & C · O & a primeira carta topográfica vetorial elaborada & productType: "SCN Carta Topográfica Vetorial"; sortField: "creationDate"; sortDirection: "ASC"; limit: 1 (opcional) \\ \hline
GO015 & C · O & cartas topográficas da mais nova para a mais antiga & productType: "SCN Carta Topográfica Matricial"; sortField: "publicationDate" ou "creationDate"; sortDirection: "DESC" \\ \hline
GO016 & C · T · O & as cartas mais recentes publicadas entre 2022 e 2023 & publicationPeriod: 2022-01-01 a 2023-12-31; sortField: "publicationDate" ou "creationDate"; sortDirection: "DESC" \\ \hline
GO017 & C · T · O & a primeira carta criada na semana passada & creationPeriod: semana passada; sortField: "creationDate"; sortDirection: "ASC"; limit: 1 (opcional) \\ \hline
GO018 & C · T · O & as 5 cartas mais antigas publicadas nos últimos 5 anos & publicationPeriod: últimos 5 anos; sortField: "publicationDate" ou "creationDate"; sortDirection: "ASC"; limit: 5 \\ \hline
GO019 & C · T · O & a última carta publicada nos últimos 6 meses & publicationPeriod: últimos 6 meses; sortField: "publicationDate"; sortDirection: "DESC"; limit: 1 (opcional) \\ \hline
\end{longtable}
\addtocounter{table}{-1}
}

\section{Consultas amb\'iguas/informais adicionais (G-A)}

Total desta fam\'ilia: \textbf{24 consultas}, listadas integralmente abaixo (a fun\c{c}\~ao geradora e o modelo de frase de cada consulta constam de \texttt{dataset.json}).

{\footnotesize
\begin{longtable}{|p{1.1cm}|p{1.3cm}|p{5.4cm}|p{6.3cm}|}
\hline
\textbf{ID} & \textbf{Cat.} & \textbf{Consulta} & \textbf{Leituras aceitas} \\ \hline
\endhead
GA001 & S · A & mapas de pe & state: "Pernambuco" \\ \hline
GA002 & S · A & mapas do pr & state: "Paraná" \\ \hline
GA003 & S · A & mapas de sc & state: "Santa Catarina" \\ \hline
GA004 & S · A & mapas de go & state: "Goiás" \\ \hline
GA005 & S · A & mapas do ce & state: "Ceará" \\ \hline
GA006 & S · A & mapas do pa & state: "Pará" \\ \hline
GA007 & S · A & mapas de goias & state: "Goiás" \\ \hline
GA008 & S · A & mapas de rondonia & state: "Rondônia" \\ \hline
GA009 & S · A & mapas do piaui & state: "Piauí" \\ \hline
GA010 & S · A & mapas do amapa & state: "Amapá" \\ \hline
GA011 & S · A & mapas do maranhao & state: "Maranhão" \\ \hline
GA012 & S · A & mapas da paraiba & state: "Paraíba" \\ \hline
GA013 & S · A & mapas de cuiaba & city: "Cuiabá" \\ \hline
GA014 & S · A & mapas de goiania & city: "Goiânia" \\ \hline
GA015 & S · A & mapas de belem & city: "Belém" \\ \hline
GA016 & S · A & mapas de brasilia & city: "Brasília" \\ \hline
GA017 & C · A & cartas do rio em 50k & state: "Rio de Janeiro"; scale: "1:50.000"; aceita-se também city no lugar de state \\ \hline
GA018 & C · A & mapas detalhados do Amazonas & scale: "1:25.000"; state: "Amazonas" \\ \hline
GA019 & S · A & cartas em pequena escala do sul & scale: "1:250.000" \\ \hline
GA020 & S · A & produtos em média escala & scale: "1:50.000" ou "1:100.000" \\ \hline
GA021 & C · A & cartas de Goiás em 25k & state: "Goiás"; scale: "1:25.000" \\ \hline
GA022 & C · A & cartas do Tocantins em 50k & state: "Tocantins"; scale: "1:50.000" \\ \hline
GA023 & C · A & cartas de Mato Grosso do Sul em 100k & state: "Mato Grosso do Sul"; scale: "1:100.000" \\ \hline
GA024 & C · A & cartas de Pernambuco em 250k & state: "Pernambuco"; scale: "1:250.000" \\ \hline
\end{longtable}
\addtocounter{table}{-1}
}

\section{Consultas amplas, com muitos produtos (G-P)}

Total desta fam\'ilia: \textbf{13 consultas}, listadas integralmente abaixo (a fun\c{c}\~ao geradora e o modelo de frase de cada consulta constam de \texttt{dataset.json}).

{\footnotesize
\begin{longtable}{|p{1.1cm}|p{1.3cm}|p{5.4cm}|p{6.3cm}|}
\hline
\textbf{ID} & \textbf{Cat.} & \textbf{Consulta} & \textbf{Leituras aceitas} \\ \hline
\endhead
GP001 & S & todas as cartas de Rondônia & state: "Rondônia" \\ \hline
GP002 & S & todas as cartas do Tocantins & state: "Tocantins" \\ \hline
GP003 & S & todas as cartas do Piauí & state: "Piauí" \\ \hline
GP004 & S & todas as cartas de Mato Grosso & state: "Mato Grosso" \\ \hline
GP005 & S & todas as cartas do Acre & state: "Acre" \\ \hline
GP006 & S & todos os produtos do 1$^\circ$ Centro de Geoinformação & supplyArea: "1$^\circ$ Centro de Geoinformação" \\ \hline
GP007 & S & todos os produtos do 2$^\circ$ Centro de Geoinformação & supplyArea: "2$^\circ$ Centro de Geoinformação" \\ \hline
GP008 & S & todos os produtos do 3$^\circ$ Centro de Geoinformação & supplyArea: "3$^\circ$ Centro de Geoinformação" \\ \hline
GP009 & S & todos os produtos do 4$^\circ$ Centro de Geoinformação & supplyArea: "4$^\circ$ Centro de Geoinformação" \\ \hline
GP010 & S & todos os produtos do 5$^\circ$ Centro de Geoinformação & supplyArea: "5$^\circ$ Centro de Geoinformação" \\ \hline
GP011 & C & cartas na escala 1:2.000 do Rio Grande do Norte & scale: "1:2.000"; state: "Rio Grande do Norte" \\ \hline
GP012 & C & cartas na escala 1:5.000 de Alagoas & scale: "1:5.000"; state: "Alagoas" \\ \hline
GP013 & C & cartas na escala 1:250.000 do Rio Grande do Sul & scale: "1:250.000"; state: "Rio Grande do Sul" \\ \hline
\end{longtable}
\addtocounter{table}{-1}
}

\section{Consultas de fronteira do dom\'inio (G-F)}

Total desta fam\'ilia: \textbf{7 consultas}, listadas integralmente abaixo (a fun\c{c}\~ao geradora e o modelo de frase de cada consulta constam de \texttt{dataset.json}).

{\footnotesize
\begin{longtable}{|p{1.1cm}|p{1.3cm}|p{5.4cm}|p{6.3cm}|}
\hline
\textbf{ID} & \textbf{Cat.} & \textbf{Consulta} & \textbf{Leituras aceitas} \\ \hline
\endhead
GF001 & F & me ajuda & não chamar a ferramenta (sem intenção de busca) \\ \hline
GF002 & F & cartas de Marte & não chamar a ferramenta (fora da Terra) \\ \hline
GF003 & F & cartas da Argentina & não chamar a ferramenta (fora do território brasileiro) \\ \hline
GF004 & F & qual a previsão do tempo para amanhã em Manaus? & não chamar a ferramenta (outro assunto) \\ \hline
GF005 & F & quanto custa uma carta topográfica? & não chamar a ferramenta (pergunta de preço, não de busca) \\ \hline
GF006 & F & mapa do tesouro pirata & não chamar a ferramenta (ficção) \\ \hline
GF007 & S & cartas do estado 42 & não chamar a ferramenta (observacional, P6: valor impossível de representar) \\ \hline
\end{longtable}
\addtocounter{table}{-1}
}

